\documentclass[11pt]{article}
\usepackage{fontspec}
\usepackage[french,english]{babel}
\usepackage{amsmath}
\usepackage{amssymb}
\usepackage{amsthm}
\usepackage{mathrsfs}
\usepackage{tikz-cd}
\usepackage{enumitem}

\newcommand{\printedpage}[1]{\marginpar{\footnotesize p.~#1}}
\newcommand{\msp}[1]{\marginpar{\footnotesize [#1]}}
\newcommand{\circled}[1]{\tikz[baseline=(c.base)]\node[draw,circle,inner sep=1pt](c){#1};}
\newcommand{\Ob}{\operatorname{Ob}}

\begin{document}

\printedpage{58}
% head: A. Grothendieck : Dérivateurs, Ch. XII, Édition des Universités de Montpellier II et Paris VII
% foot: Version 7 décembre 2002

is \textit{surjective.} (E.g. if $S'=\emptyset$, we have $X'=\emptyset$, hence $f'\in W$, without that implying $f\in W$!) The fact that in the cocartesian case of W(7 bis), $i_0\in W\Longrightarrow i\in W$ must be regarded as ``going without saying'' - otherwise the square does not deserve, from the $W$-homotopic point of view, the name ``cocartesian''. On the other hand, the reverse implication $i\in W\Longrightarrow i_0\in W$ appears to me as a very particular feature of the derivator HOT$_W$ (when it has it), and notably % [sic] dérivatuer -> dérivateur, misprint
of HOT $=$ HOT$_{W_\infty}$. The beginning
\msp{page 80}
of the remark simply means that the dual derivator HOT$^\circ$ does not have this same property. I must refer in this connection to the two axioms for derivators corresponding to W(7 bis), W(7), which I will probably do shortly, so as to give a more precise meaning to the present comments.

\subsection*{2.2 Application to cocartesian squares deduced from an open or closed immersion.}

The axioms W(7), W(7 bis) (of the cube, and of the cocartesian square) concern cocartesian squares of a very particular nature. We are now going to see how to deduce from them statements of the same type for much more general cocartesian squares, including for squares which have nothing cocartesian about them (in Cat) - and which we will call % [sic] appelerons -> appellerons, misspelling
``homotopically cocartesian'' (relative to $W$), or simply $W$-cocartesian. (There will also be $W$-cartesian squares one day \dots) In the case where $W=W_\infty$, I will also call them Hot-cocartesian. % [sic] appelerai -> appellerai, misspelling

I will begin with the particular case of a cocartesian square associated with an open immersion (diagram ($*$), p. 74):
\[
(*)\qquad\qquad
\begin{tikzcd}
U \arrow[r, hook, "i"] \arrow[d, "p"'] & X \arrow[d, "q"'] & T \arrow[l, hook', dotted, "j"'] \arrow[d, dotted] \\
U' \arrow[r, hook, "i'"] & X' & T'\simeq T. \arrow[l, hook', dotted, "j'"']
\end{tikzcd}
\]
\msp{page 81}

The big trap here that we have to manage to exorcise is that it can happen that $i\in W$, without our having $i'\in W$, and also $p\in W$ without our having $q\in W$. I had convinced myself of this in 1983 while writing Pursuing Stacks (during the weeks of ``scratchwork''), and don't feel like looking for a counterexample here, unless it presents itself on its own (which it probably will).

I am going to give an important case where these troubles disappear:

\textbf{Proposition 6:} \textit{Consider the closures $\bar U$ and $\overline{U'}$ of $U$ and $U'$ in $X$, $X'$ respectively, and suppose that the inclusions}
\[
U\xrightarrow{i_0}\bar U,\qquad U'\xrightarrow{i'_0}\overline{U'}
\]
\printedpage{59}
% head: A. Grothendieck : Dérivateurs, Ch. XII, Édition des Universités de Montpellier II et Paris VII
% foot: Version 7 décembre 2002
\textit{are in $W$} (\footnote[64]{cf. th. 3 p. 107 for general conditions ensuring that this hypothesis is satisfied.}). \textit{Then}
\[
\left\{
\begin{array}{llll}
\circled{\textit{a}}\ i\in W & \iff & i'\in W & \textit{if we have }\mathrm{W}(7\text{ bis})\\
\circled{\textit{b}}\ p\in W & \iff & q\in W & \textit{if we have }\mathrm{W}(7).
\end{array}
\right.
\]

\textsc{Proof}. Let us first consider case \circled{a} of $i$, $i'$. Let
\[
k:\bar U\hookrightarrow X
\]
be the inclusion. The triangle
\begin{tikzcd}
U \arrow[r, "i_0"] \arrow[rr, bend right=8, "i"'] & \bar U \arrow[r, "k"] & X
\end{tikzcd}
, where $i_0\in W$, shows that
\[
i\in W\iff k\in W.
\]

On the other hand, $X$ is the union of the closed subcategories $\bar U$, $T$, so we have a cocartesian square in Cat
\[
(*)\qquad\qquad
\begin{tikzcd}
\bar U\cap T=\dot T \arrow[r, "k_0"] \arrow[d, "j_0"'] & T \arrow[d, "j"] \\
\bar U \arrow[r, "k"] & X.
\end{tikzcd}
\]

The axiom W(7 bis) (in its dual form (i.e. applied to the square formed by the opposite categories \dots)) implies
\msp{page 82}
that
\[
k\in W\iff k_0\in W,
\]
whence
\[
i\in W\iff (k_0:\dot T\longrightarrow T)\in W.
\]
Likewise we find, applying this to $(X',U')$ in place of $(X,U)$ % [sic] appliqunt -> appliquant, misprint
\[
i'\in W\iff (k'_0:\dot T'\longrightarrow T')\in W.
\]
Now we have $T\xrightarrow{\sim}T'$ (isomorphism), and by proposition 4 this isomorphism takes $\dot T$ to $\dot T'$. Thus
\[
k_0\in W\iff k'_0\in W\qquad\quad\text{(trivially)},
\]
whence a).

Let us prove b). The commutative square induced by ($*$), p. 80,
\printedpage{60}
% head: A. Grothendieck : Dérivateurs, Ch. XII, Édition des Universités de Montpellier II et Paris VII
% foot: Version 7 décembre 2002
\[
\begin{tikzcd}
U \arrow[r, hook, "i_0"] \arrow[d, "p"'] & \bar U \arrow[d, "\bar p"] \\
U' \arrow[r, hook, "i'_0"] & \overline{U'},
\end{tikzcd}
\]
where $i_0,i'_0\in W$ by hypothesis, shows that
\[
p\in W\iff\bar p\in W,
\]
so everything comes down to proving that
\[
\bar p\in W\overset{?}{\iff}q\in W.
\]
But $q$ induces a homomorphism from the square ($*$) deduced from $(X,\bar U,T)$ (with $X=\bar U\cup T$), into the similar square for $X'$
\[
\begin{tikzcd}
\overline{U'}\cap T'=\dot T' \arrow[r, "k'_0"] \arrow[d, "j'_0"'] & T \arrow[d, "j'"] \\ % [sic] T -> T', the square for X' should have the corner T' (although T'\simeq T)
\overline{U'} \arrow[r, "k'"] & X',
\end{tikzcd}
\]
the vertical arrows being denoted by $q_{\dot T}$, $q_T$, $\bar p$, $q$.
\msp{page 83}
We know that $q_{\dot T}$, $q_T$ are isomorphisms, hence in $W$, so the ``five lemma'' W(7) assures us that $\bar p\in W\iff q\in W$, q.e.d.

\textbf{Comment.} This proposition leads us to look for conditions on an open subcategory $U$ of an $X$, which ensure that the inclusion $i_0:U\longrightarrow\bar U$ is \texttt{[in]} $W$, and \textit{that this property is preserved by every cobase change $U\longrightarrow U'$.} Moreover, we seek conditions which are of a local nature on $\bar U$, more precisely on $\bar U\cap T=\dot U=\dot T$. We can motivate this concern (a posteriori) by saying that we would like the property $i_0\in W$ to be preserved, not only by an arbitrary cobase change $U\longrightarrow U'$, but also, in addition, by passing from $X$ to any full subcategory $X_1$ containing $U$, thus giving rise to a situation
\[
U\overset{i_Y}{\hookrightarrow}X_1,\qquad U\xrightarrow{(i_Y)_0}\bar U^{X_1}=\bar U\cap X_1
\]
\texttt{[}$Y=X_1$\texttt{?]}, similar to the starting situation. More generally still, we would like it to be preserved on passing from the triple $(X,U,T)$ to an ``induced'' triple $(X_1,U,T_1)$, by an arbitrary arrow $g:T_1\longrightarrow T$ with target $T$ in Cat, giving rise to the pair of arrows
\[
U\xrightarrow{\mathrm{id}}U,\ g:T_1\longrightarrow T
\]
(cf. pages 60, 61). The preceding case is the one where
\printedpage{61}
% head: A. Grothendieck : Dérivateurs, Ch. XII, Édition des Universités de Montpellier II et Paris VII
% foot: Version 7 décembre 2002
\msp{page 84}
$T_1$ is a full subcategory of $T$, $g$ the inclusion arrow, Let us note in this connection the sorital proposition % [sic] inclusion, Notons -> inclusion. Notons, comma before a capital letter

\textbf{Proposition 7.} \textit{Equivalent conditions on the open subcategory $U$ of an object $X$ of} Cat (assuming W(7 bis)):

\begin{enumerate}[label={}, leftmargin=*]
\item[\textit{a)}] \textit{The inclusion $i_0:U\longrightarrow\bar U$ is $W$-aspheric, i.e. for every $t\in\dot U=\bar U\backslash U$, the category $U/t$ is $W$-aspheric.}
\item[\textit{b)}] \textit{The inclusion $i_0:U\longrightarrow\bar U$ is in $W$, and this condition remains satisfied, when we replace $(X,U)$ by $(X_1,U)$, where $(X_1,U)$ is deduced from $(X,U)$ by means of an arrow $T_1\longrightarrow T(=X\backslash U)$ in} Cat\textit{, as recalled above.}
\item[\textit{b$'$)}] \textit{As in b), in the case where $T_1$ is the one-point category $e$, and where moreover its image is in $\dot U$.}
\end{enumerate}

\textsc{Proof}. It is obvious that if in b), b$'$) we replace the condition ``$i_0\in W$'' by the stronger condition ``$i_0$ is $W$-aspheric'', the four \texttt{[rather three]} conditions a) to b$'$) are equivalent. It remains to prove that the condition that the $U/t$ be $W$-aspheric ($t\in\bar U$ - but this is trivial for $t\in U$, so it suffices to express it for $t\in\bar U\backslash U=\dot U$) is already implied by the weakest b$'$ and we guess that we have this:

\msp{page 85}

\textbf{Corollary.} (Assuming W(7 bis).) \textit{Let $X$ be in} Cat\textit{, $U$ an open subcategory of $X$ such that $X\backslash U=\overline{\{t\}}$ is the final category. Then the inclusion $U\xrightarrow{i}X$ is in $W$ if and only if it is $W$-aspheric, i.e. if and only if $U/t$ is $W$-aspheric (where $t$ is the unique element of $T$).}

We prove this corollary later (\footnote[65]{cf. cor. 4 of proposition 15, page 149.}). Proposition 7 will be of no use to us until then - it will have only a purely aesthetic role to play, to justify our deliberate intention of looking for conditions on $i:U\longrightarrow X$ (to ensure that $U\longrightarrow\bar U$ is in $W$ and remains so, when we replace $(X,U)$ by other pairs $(X_1,U_1)$ which are deduced from it by certain operations \dots) which
\msp{page 86}
are of a \textit{local nature} on $\bar U$ (in fact, at the points of $\dot U$), and which in any event imply the \textit{$W$-asphericity} of $U\longrightarrow\bar U$.

I had hoped that this asphericity condition (suggested by the notion ``of a Dwyer immersion'') would be stable under an arbitrary cobase change $U\longrightarrow U'$, which would at the same time have ensured that after such an operation, the open immersion $i':U'\longrightarrow X'$ obtained still satisfies the condition that $i'_0:U'\longrightarrow\overline{U'}$ is in $W$. We will see that this expectation is unfounded - the condition of $W$-asphericity is \textit{not} stable under cobase change, and it does \textit{not} even suffice to ensure that after a cobase change, $i'_0:U'\longrightarrow\overline{U'}$ is
\printedpage{62}
% head: A. Grothendieck : Dérivateurs, Ch. XII, Édition des Universités de Montpellier II et Paris VII
% foot: Version 7 décembre 2002
in $W$. A stronger notion will be needed, which is suggested by the case where $T=\{t\}$ is the final category and contained in $\bar U$ (the case to which we tried to reduce ourselves in proposition 7). In this case, as we said (preceding corollary), the condition that $i:U\longrightarrow\bar U=X$ is in $W$, is equivalent to the $W$-asphericity of $i$, i.e. to the $W$-asphericity of $U/t$, and likewise for $i':U'\longrightarrow\overline{U'}=X'$: it is in $W$ if and only if $U'/t'$ is $W$-aspheric, where
\msp{page 87}
$t'$ is the unique element of $T'(\simeq T)=X'\backslash U'$. So we need conditions on $U/t$ which ensure that not only $U/t$, but also all the similar categories $U'/t$ \texttt{[rather} $U'/t'$\texttt{]}, are $W$-aspheric. In the general case where $T$ is arbitrary, it is this condition that we are going to impose on the $U/t$ ($t\in\Ob T$), and it will turn out that this time, they are stable \texttt{[rather it is stable]} under every cobase change $U\xrightarrow{p}U'$, i.e. that the $U'/t'$ ($t'\in\Ob T'$) will then satisfy the same condition. It is this condition that we are now going to reformulate.

\subsection*{2.3 Totally $W$-aspheric objects of \textnormal{Cat}, and totally $W$-aspheric morphisms.}

Let us take up again the situation
\[
U\overset{i}{\hookrightarrow}X\overset{j}{\hookleftarrow}T,\qquad\qquad\text{with }T=\{t\}\text{, final category, }t\in\bar U\text{, i.e. }X=\bar U,
\]
and that of a cobase change $U\longrightarrow U'$, giving rise to a cocartesian square in Cat
\[
(*)\qquad\qquad
\begin{tikzcd}
U \arrow[r, hook, "i"] \arrow[d, "p"'] & X \arrow[d, "q"'] & T=\{t\} \arrow[l, hook', dotted, "j"'] \arrow[d, dotted, "\mathrm{iso}"] \\
U' \arrow[r, hook, "i'"] & X' & T'=\{t'\}. \arrow[l, hook', dotted, "j'"']
\end{tikzcd}
\]
The gluing datum $H$ between $U$ and $T$, $H:\underbrace{U^\circ\times T}_{\simeq U^\circ}\longrightarrow\mathrm{Ens}$, can be interpreted as a presheaf on $U$, $H\in\Ob U^\wedge:U^\circ\longrightarrow\mathrm{Ens}$,
\msp{page 88}
and we have
\[
U/t\simeq U/H,
\]
where $U/H$ is interpreted in terms of the canonical embedding
\[
U\hookrightarrow U^\wedge.
\]
Likewise, the gluing between $U'$ and $T'$ will be given by a presheaf
\[
H'\in\Ob U^{\prime\wedge}:U^{\prime\circ}\longrightarrow\mathrm{Ens},
\]

\end{document}
